
\documentclass{article}
\usepackage{colortbl}
\usepackage{makecell}
\usepackage{multirow}
\usepackage{supertabular}

\begin{document}

\newcounter{utterance}

\centering \large Interaction Transcript for game `dond', experiment `semi\_en', episode 10 with Qwen3.5{-}4B.
\vspace{24pt}

{ \footnotesize  \setcounter{utterance}{1}
\setlength{\tabcolsep}{0pt}
\begin{supertabular}{c@{$\;$}|p{.15\linewidth}@{}p{.15\linewidth}p{.15\linewidth}p{.15\linewidth}p{.15\linewidth}p{.15\linewidth}}
    \# & $\;$A & \multicolumn{4}{c}{Game Master} & $\;\:$B\\
    \hline

    \theutterance \stepcounter{utterance}  
    & & \multicolumn{4}{p{0.6\linewidth}}{
        \cellcolor[rgb]{0.9,0.9,0.9}{
            \makecell[{{p{\linewidth}}}]{
                \texttt{\tiny{[P1$\langle$GM]}}
                \texttt{You are playing a negotiation game in which you have to agree on how to divide a set of items among you and another player.} \\
\\ 
\texttt{Rules:} \\
\texttt{(a) You and the other player are given a set of items. Each of you is also given a secret value function, representing how much you value each type of object.} \\
\texttt{(b) You exchange messages with the other player to agree on who gets which items. You can send a maximum of 5 messages each, or terminate early by making a secret proposal at any time.} \\
\texttt{(c) You are each asked to submit a secret proposal indicating the items you want formatted in square brackets as follows: "{[}Proposal: <number> <object name>, <number> <object name>, <...>{]}"} \\
\texttt{(d) If your proposals are complementary, i.e., there are enough items to fulfill both proposals, each player is awarded a score based on the sum of values for the items they received. Otherwise, both of you get zero points.} \\
\texttt{(e) Your goal is to maximize the score you receive.} \\
\\ 
\texttt{Let us start.} \\
\\ 
\texttt{The set of available items is:} \\
\\ 
\texttt{2 remotes, 2 magnets, 4 watches} \\
\\ 
\texttt{Your secret value function is:} \\
\\ 
\texttt{remote: 3, magnet: 2, watch: 0} \\
\\ 
\texttt{IMPORTANT: Your messages, unless it is the secret proposal, are directly transmitted to the other player, so do not include any response to the rules or text announcing your message. To make a secret proposal, use the indicated format. Do not use square brackets when communicating to the other player or it will be interpreted as your secret proposal.} \\
            }
        }
    }
    & \\ \\

    \theutterance \stepcounter{utterance}  
    & \multicolumn{4}{p{0.6\linewidth}}{
        \cellcolor[rgb]{0.9,0.9,0.9}{
            \makecell[{{p{\linewidth}}}]{
                \texttt{\tiny{[P1$\rangle$GM]}}
                \texttt{I will propose a division that maximizes my value while ensuring the other player also receives a positive value (to avoid the zero{-}point penalty).} \\
\\ 
\texttt{My valuation:} \\
\texttt{{-} Remote: 3} \\
\texttt{{-} Magnet: 2} \\
\texttt{{-} Watch: 0} \\
\\ 
\texttt{Total value of all items: 2*3 + 2*2 + 4*0 = 6 + 4 = 10.} \\
\\ 
\texttt{If I take all 2 remotes and 2 magnets, I get 6 + 4 = 10 points. The other player gets 4 watches, which are worth 0 to them. This might lead to a disagreement or a zero score if they reject it.} \\
\\ 
\texttt{A better strategy is to split the high{-}value items.} \\
\texttt{Option A: I take 1 remote (3 pts) and 1 magnet (2 pts) = 5 pts. Other player takes 1 remote (3 pts) and 1 magnet (2 pts) = 5 pts. The 4 watches are left over (worth 0 to both). Total score = 5 + 5 = 10. This is a fair split.} \\
\\ 
\texttt{Let's try to give the other player something they might value, but since I don't know their values, I must assume they value the items at least somewhat. However, the rule says "If your proposals are complementary... awarded a score... Otherwise... zero points." This implies if we don't agree on a partition that covers all items (or at least a valid partition), we get zero. Actually, re{-}reading (d): "If your proposals are complementary, i.e., there are enough items to fulfill both proposals". This usually means if my proposal is a set of items and your proposal is a disjoint set of items, and the union is the whole set (or at least valid), we succeed.} \\
\\ 
\texttt{Let's propose a specific split.} \\
\texttt{I will propose: 1 remote, 1 magnet.} \\
\texttt{I hope the other player proposes: 1 remote, 1 magnet.} \\
\texttt{Then the 4 watches are ignored or given to someone? The rule says "divide a set of items". Usually, in these games, you propose a partition.} \\
\texttt{Let's propose: 1 remote, 1 magnet.} \\
\texttt{If the other player proposes: 1 remote, 1 magnet, and maybe some watches?} \\
\texttt{Wait, if I propose "1 remote, 1 magnet", and they propose "1 remote, 1 magnet", we have used 2 remotes and 2 magnets. The 4 watches are unused. Is that allowed? "Divide a set of items" implies all items should be distributed.} \\
\texttt{Let's try to include the watches to make it a full partition.} \\
\texttt{Since watches are worth 0 to me, I can give them to the other player to make the deal work, or take them if they are valuable to them.} \\
\texttt{But I don't} \\
            }
        }
    }
    & & \\ \\

    \theutterance \stepcounter{utterance}  
    & & & \multicolumn{2}{p{0.3\linewidth}}{
        \cellcolor[rgb]{0.9,0.9,0.9}{
            \makecell[{{p{\linewidth}}}]{
                \texttt{\tiny{[GM$|$GM]}}
                \texttt{continue} \\
            }
        }
    }
    & & \\ \\

    \theutterance \stepcounter{utterance}  
    & & & \multicolumn{2}{p{0.3\linewidth}}{
        \cellcolor[rgb]{0.9,0.9,0.9}{
            \makecell[{{p{\linewidth}}}]{
                \texttt{\tiny{[GM$|$GM]}}
                \texttt{I will propose a division that maximizes my value while ensuring the other player also receives a positive value (to avoid the zero{-}point penalty).} \\
\\ 
\texttt{My valuation:} \\
\texttt{{-} Remote: 3} \\
\texttt{{-} Magnet: 2} \\
\texttt{{-} Watch: 0} \\
\\ 
\texttt{Total value of all items: 2*3 + 2*2 + 4*0 = 6 + 4 = 10.} \\
\\ 
\texttt{If I take all 2 remotes and 2 magnets, I get 6 + 4 = 10 points. The other player gets 4 watches, which are worth 0 to them. This might lead to a disagreement or a zero score if they reject it.} \\
\\ 
\texttt{A better strategy is to split the high{-}value items.} \\
\texttt{Option A: I take 1 remote (3 pts) and 1 magnet (2 pts) = 5 pts. Other player takes 1 remote (3 pts) and 1 magnet (2 pts) = 5 pts. The 4 watches are left over (worth 0 to both). Total score = 5 + 5 = 10. This is a fair split.} \\
\\ 
\texttt{Let's try to give the other player something they might value, but since I don't know their values, I must assume they value the items at least somewhat. However, the rule says "If your proposals are complementary... awarded a score... Otherwise... zero points." This implies if we don't agree on a partition that covers all items (or at least a valid partition), we get zero. Actually, re{-}reading (d): "If your proposals are complementary, i.e., there are enough items to fulfill both proposals". This usually means if my proposal is a set of items and your proposal is a disjoint set of items, and the union is the whole set (or at least valid), we succeed.} \\
\\ 
\texttt{Let's propose a specific split.} \\
\texttt{I will propose: 1 remote, 1 magnet.} \\
\texttt{I hope the other player proposes: 1 remote, 1 magnet.} \\
\texttt{Then the 4 watches are ignored or given to someone? The rule says "divide a set of items". Usually, in these games, you propose a partition.} \\
\texttt{Let's propose: 1 remote, 1 magnet.} \\
\texttt{If the other player proposes: 1 remote, 1 magnet, and maybe some watches?} \\
\texttt{Wait, if I propose "1 remote, 1 magnet", and they propose "1 remote, 1 magnet", we have used 2 remotes and 2 magnets. The 4 watches are unused. Is that allowed? "Divide a set of items" implies all items should be distributed.} \\
\texttt{Let's try to include the watches to make it a full partition.} \\
\texttt{Since watches are worth 0 to me, I can give them to the other player to make the deal work, or take them if they are valuable to them.} \\
\texttt{But I don't} \\
            }
        }
    }
    & & \\ \\

    \theutterance \stepcounter{utterance}  
    & & \multicolumn{4}{p{0.6\linewidth}}{
        \cellcolor[rgb]{0.9,0.9,0.9}{
            \makecell[{{p{\linewidth}}}]{
                \texttt{\tiny{[GM$\rangle$P2]}}
                \texttt{You are playing a negotiation game in which you have to agree on how to divide a set of items among you and another player.} \\
\\ 
\texttt{Rules:} \\
\texttt{(a) You and the other player are given a set of items. Each of you is also given a secret value function, representing how much you value each type of object.} \\
\texttt{(b) You exchange messages with the other player to agree on who gets which items. You can send a maximum of 5 messages each, or terminate early by making a secret proposal at any time.} \\
\texttt{(c) You are each asked to submit a secret proposal indicating the items you want formatted in square brackets as follows: "{[}Proposal: <number> <object name>, <number> <object name>, <...>{]}"} \\
\texttt{(d) If your proposals are complementary, i.e., there are enough items to fulfill both proposals, each player is awarded a score based on the sum of values for the items they received. Otherwise, both of you get zero points.} \\
\texttt{(e) Your goal is to maximize the score you receive.} \\
\\ 
\texttt{Let us start.} \\
\\ 
\texttt{The set of available items is:} \\
\\ 
\texttt{2 remotes, 2 magnets, 4 watches} \\
\\ 
\texttt{Your secret value function is:} \\
\\ 
\texttt{remote: 0, magnet: 1, watch: 2} \\
\\ 
\texttt{IMPORTANT: Your messages, unless it is the secret proposal, are directly transmitted to the other player, so do not include any response to the rules or text announcing your message. To make a secret proposal, use the indicated format. Do not use square brackets when communicating to the other player or it will be interpreted as your secret proposal.} \\
\\ 
\\ 
\texttt{I will propose a division that maximizes my value while ensuring the other player also receives a positive value (to avoid the zero{-}point penalty).} \\
\\ 
\texttt{My valuation:} \\
\texttt{{-} Remote: 3} \\
\texttt{{-} Magnet: 2} \\
\texttt{{-} Watch: 0} \\
\\ 
\texttt{Total value of all items: 2*3 + 2*2 + 4*0 = 6 + 4 = 10.} \\
\\ 
\texttt{If I take all 2 remotes and 2 magnets, I get 6 + 4 = 10 points. The other player gets 4 watches, which are worth 0 to them. This might lead to a disagreement or a zero score if they reject it.} \\
\\ 
\texttt{A better strategy is to split the high{-}value items.} \\
\texttt{Option A: I take 1 remote (3 pts) and 1 magnet (2 pts) = 5 pts. Other player takes 1 remote (3 pts) and 1 magnet (2 pts) = 5 pts. The 4 watches are left over (worth 0 to both). Total score = 5 + 5 = 10. This is a fair split.} \\
\\ 
\texttt{Let's try to give the other player something they might value, but since I don't know their values, I must assume they value the items at least somewhat. However, the rule says "If your proposals are complementary... awarded a score... Otherwise... zero points." This implies if we don't agree on a partition that covers all items (or at least a valid partition), we get zero. Actually, re{-}reading (d): "If your proposals are complementary, i.e., there are enough items to fulfill both proposals". This usually means if my proposal is a set of items and your proposal is a disjoint set of items, and the union is the whole set (or at least valid), we succeed.} \\
\\ 
\texttt{Let's propose a specific split.} \\
\texttt{I will propose: 1 remote, 1 magnet.} \\
\texttt{I hope the other player proposes: 1 remote, 1 magnet.} \\
\texttt{Then the 4 watches are ignored or given to someone? The rule says "divide a set of items". Usually, in these games, you propose a partition.} \\
\texttt{Let's propose: 1 remote, 1 magnet.} \\
\texttt{If the other player proposes: 1 remote, 1 magnet, and maybe some watches?} \\
\texttt{Wait, if I propose "1 remote, 1 magnet", and they propose "1 remote, 1 magnet", we have used 2 remotes and 2 magnets. The 4 watches are unused. Is that allowed? "Divide a set of items" implies all items should be distributed.} \\
\texttt{Let's try to include the watches to make it a full partition.} \\
\texttt{Since watches are worth 0 to me, I can give them to the other player to make the deal work, or take them if they are valuable to them.} \\
\texttt{But I don't} \\
            }
        }
    }
    & \\ \\

    \theutterance \stepcounter{utterance}  
    & & & \multicolumn{4}{p{0.6\linewidth}}{
        \cellcolor[rgb]{0.9,0.9,0.9}{
            \makecell[{{p{\linewidth}}}]{
                \texttt{\tiny{[GM$\langle$P2]}}
                \texttt{{[}Proposal: 1 remote, 1 magnet{]}} \\
            }
        }
    }
     \\ \\

    \theutterance \stepcounter{utterance}  
    & & & \multicolumn{2}{p{0.3\linewidth}}{
        \cellcolor[rgb]{0.9,0.9,0.9}{
            \makecell[{{p{\linewidth}}}]{
                \texttt{\tiny{[GM$|$GM]}}
                \texttt{proposal} \\
            }
        }
    }
    & & \\ \\

    \theutterance \stepcounter{utterance}  
    & & & \multicolumn{2}{p{0.3\linewidth}}{
        \cellcolor[rgb]{0.9,0.9,0.9}{
            \makecell[{{p{\linewidth}}}]{
                \texttt{\tiny{[GM$|$GM]}}
                [1, 1, 0]
            }
        }
    }
    & & \\ \\

    \theutterance \stepcounter{utterance}  
    & & \multicolumn{4}{p{0.6\linewidth}}{
        \cellcolor[rgb]{0.9,0.9,0.9}{
            \makecell[{{p{\linewidth}}}]{
                \texttt{\tiny{[P1$\langle$GM]}}
                \texttt{The other player has submitted a secret proposal. It is now your turn to submit a proposal of your own.} \\
            }
        }
    }
    & \\ \\

    \theutterance \stepcounter{utterance}  
    & \multicolumn{4}{p{0.6\linewidth}}{
        \cellcolor[rgb]{0.9,0.9,0.9}{
            \makecell[{{p{\linewidth}}}]{
                \texttt{\tiny{[P1$\rangle$GM]}}
                \texttt{{[}Proposal: 1 remote, 1 magnet{]}} \\
            }
        }
    }
    & & \\ \\

    \theutterance \stepcounter{utterance}  
    & & & \multicolumn{2}{p{0.3\linewidth}}{
        \cellcolor[rgb]{0.9,0.9,0.9}{
            \makecell[{{p{\linewidth}}}]{
                \texttt{\tiny{[GM$|$GM]}}
                \texttt{proposal} \\
            }
        }
    }
    & & \\ \\

    \theutterance \stepcounter{utterance}  
    & & & \multicolumn{2}{p{0.3\linewidth}}{
        \cellcolor[rgb]{0.9,0.9,0.9}{
            \makecell[{{p{\linewidth}}}]{
                \texttt{\tiny{[GM$|$GM]}}
                [1, 1, 0]
            }
        }
    }
    & & \\ \\

    \theutterance \stepcounter{utterance}  
    & & & \multicolumn{2}{p{0.3\linewidth}}{
        \cellcolor[rgb]{0.9,0.9,0.9}{
            \makecell[{{p{\linewidth}}}]{
                \texttt{\tiny{[GM$|$GM]}}
                [[1, 1, 0], [1, 1, 0]]
            }
        }
    }
    & & \\ \\

\end{supertabular}
}

\end{document}
